\section{Solution}
\label{sec:solution}

The controller is a deterministic, model-based planner. At each rally it
\emph{locks} a swing plan as soon as the ball begins to head toward it,
and then executes that plan open-loop in joint space. The decision to
plan once and lock --- rather than continuously re-plan --- is deliberate:
early experimentation showed that re-solving inverse kinematics on every
control step caused the arm to chatter between two close IK branches and
miss the ball entirely.

Figure~\ref{fig:pipeline} summarises the pipeline; the rest of this
section explains each block.

\begin{figure}[H]
\centering
\begin{tikzpicture}[
  block/.style={rectangle, draw=darkRed!85, fill=darkRed!8,
                rounded corners=3pt, minimum height=10mm, minimum width=22mm,
                align=center, font=\small\sffamily},
  decision/.style={diamond, draw=darkRed!85, fill=darkRed!8,
                   inner sep=1pt, align=center, font=\small\sffamily,
                   aspect=2},
  arrow/.style={-stealth, thick, darkRed!85},
  node distance=8mm and 14mm,
]
  \node[block] (obs)   {Observation\\$\bball,\vball,\jq$};
  \node[decision, right=of obs] (head) {coming\\to us?};
  \node[block, right=of head]   (plan) {Build swing plan\\(\S\ref{sub:plan})};
  \node[block, below=of plan]   (sched){Schedule swing\\(\S\ref{sub:sched})};
  \node[block, left=of sched]   (exec) {Execute joint\\target};
  \node[block, left=of exec]    (ready){Hold ready pose};

  \draw[arrow] (obs)   -- (head);
  \draw[arrow] (head)  -- node[above, font=\scriptsize]{yes} (plan);
  \draw[arrow] (head)  |- node[left, font=\scriptsize, pos=0.7]{no} (ready);
  \draw[arrow] (plan)  -- (sched);
  \draw[arrow] (sched) -- (exec);
  \draw[arrow] (ready) -- (exec);
  \draw[arrow] (exec)  -- ++(0,-1.0) -| (obs);
\end{tikzpicture}
\caption{Per-step pipeline of the \texttt{RobotPlayer} controller.
The plan is built at most once per rally; subsequent steps only execute
the cached joint targets.}
\label{fig:pipeline}
\end{figure}

\subsection{Ball-trajectory prediction}
\label{sub:predict}

Between contacts the ball obeys~\eqref{eq:gravity}. Given current state
$(\bball_{0},\vball_{0})$ the trajectory is
\begin{equation}
  \bball(t)=\bball_{0}+\vball_{0}\,t-\tfrac{1}{2}g\,t^{2}\mathbf{e}_{z},\qquad
  \vball(t)=\vball_{0}-g\,t\,\mathbf{e}_{z}.
  \label{eq:projectile}
\end{equation}
The first table bounce occurs at the smallest positive root of
$b_{z}(t)=z_{T}$, i.e.
\begin{equation}
  \tfrac{1}{2}g\,t_{B}^{2}-v_{0z}\,t_{B}+(b_{0z}-z_{T})=0
  \;\Rightarrow\;
  t_{B}=\frac{v_{0z}+\sqrt{v_{0z}^{2}-2g(b_{0z}-z_{T})}}{g}.
  \label{eq:bounce}
\end{equation}
The bounce location is
$\mathbf{b}_{B}=\bball(t_{B})$ and the pre-bounce vertical
velocity is $v_{Bz}^{-}=v_{0z}-g\,t_{B}$. We model the
table--ball contact as Newtonian restitution with coefficient
$e_{B}=0.87$ (matched to the simulator within tuning tolerance), so the
post-bounce vertical velocity is $v_{Bz}^{+}=e_{B}|v_{Bz}^{-}|$ while the
horizontal velocity is preserved up to a small empirical damping
$\eta_{V}=0.82$:
\begin{equation}
  \vball^{+}=\bigl(\eta_{V}v_{0x},\;v_{0y},\;e_{B}|v_{Bz}^{-}|\bigr).
  \label{eq:postbounce}
\end{equation}
The damping factor $\eta_{V}$ corrects for PyBullet's small linear
damping on the ball, which causes the realised $|v_{x}|$ to decay
roughly $15$--$18\%$ over a typical $0.7$\,s flight.

From the bounce location we propagate~\eqref{eq:projectile} forward to
the agent's \emph{strike plane}
\[
  x_{s}=b_{\mathrm{base}}^{x}+(-s)\,d_{\mathrm{strike}},
  \qquad s=\begin{cases}-1 & \text{Player 0}\\+1 & \text{Player 1}\end{cases},
\]
to find the intercept time $t_{s}$ and the predicted strike point
$\mathbf{p}_{s}=(x_{s},y_{s},z_{s})$. The strike $z$ is clamped to a
safe envelope $[z^{\min},z^{\max}]=[0.96,1.20]$\,m to keep the paddle
away from the table and within reach.

\subsection{Strike pose and swing geometry}
\label{sub:plan}

The paddle face must point toward the opponent and be tilted upward by an
angle $\phi$ to help the ball clear the net. In the world frame this
sets a target paddle orientation
\begin{equation}
  \mathbf{n}_{p}=\bigl(-s\cos\phi,\;0,\;\sin\phi\bigr),
\end{equation}
where $s=\pm 1$ is the side sign defined above and
$\phi=0.65$\,rad ($\approx 37^{\circ}$) was selected by the parameter
sweep of Section~\ref{sub:tune}. The desired paddle EE pose at the
strike point is therefore $(\mathbf{p}_{s},\,\mathbf{n}_{p})$.

To deliver momentum to the ball we additionally specify two staging
poses around the strike point:
\begin{align}
  \text{windup}    \quad &\mathbf{p}_{w}=\mathbf{p}_{s}+(s\,d_{W},0,-0.03), \\
  \text{follow-through}\quad &\mathbf{p}_{f}=\mathbf{p}_{s}-(s\,d_{F},0,-d_{L}),
\end{align}
with $d_{W}=0.05$, $d_{F}=0.10$ and lift $d_{L}=0.12$\,m. The arm holds
$\mathbf{p}_{w}$ behind the strike point, then sweeps forward through
$\mathbf{p}_{s}$ to $\mathbf{p}_{f}$. Because the paddle is rigidly
attached, this sweep imparts paddle velocity at contact (in the
$+s$-toward-opponent direction and slightly upward), producing a
returnable shot without needing impedance or torque control.

\subsection{Inverse kinematics with branch consistency}
\label{sub:ik}

The paddle pose is a six-dimensional task on a seven-joint arm, so the
manipulator is kinematically redundant. PyBullet's
\texttt{calculateInverseKinematics} uses damped-least-squares with a
rest-pose null-space bias~\cite{buss2004ik}; given a target
$(\mathbf{p},\mathbf{R})$ it returns a joint vector $\jq^{\star}$ close
to the rest pose. Two practical issues had to be solved.

\paragraph{Branch consistency: the mirror-IK trick.}
For the right-hand robot (base yaw $\pi$ about $z$), PyBullet's IK
seeded from the live joint state converges to a contorted joint branch
in which the wrist joint $\theta_{7}$ sits near its limit
$+\pi$~rad. Although the end-effector position and orientation match
the request \emph{at the strike point}, the linear interpolation in
joint space between the windup branch and the follow-through branch
produces a Cartesian path with a $\sim 10$\,cm side-to-side excursion in
$y$; the paddle misses the ball entirely.

The two robot bases are placed symmetrically about the origin, so the
remedy is purely geometric. Both robots are kinematically identical:
applying the same joint vector $\jq$ to the unrotated base
$\mathbf{R}_{1}=\mathbf{I}$ and to the rotated base
$\mathbf{R}_{2}=\eR{z}(\pi)$ produces end-effector poses related by
$\eR{z}(\pi)$. We therefore always solve IK on Robot~1's kinematic
chain. For Player~1 (Robot~2) we first mirror the world target through
the origin,
\begin{equation}
  (x,y,z)\longmapsto(-x,-y,z),
  \label{eq:mirror}
\end{equation}
run IK against Robot~1 with the mirrored target, and apply the resulting
joint vector to Robot~2. The composition
$\eR{z}(\pi)\cdot\Phi_{1}(\jq)$ recovers the desired mirror world pose
without ever asking PyBullet to solve a yaw-$\pi$ IK. As shown later
(Section~\ref{sub:results_p1}), this single change raises Player~1's
paddle-contact rate from $11.8\%$ to $86.3\%$.

\paragraph{Snapshot/restore.}
Because PyBullet's IK is iterative and seeded from the robot's
\emph{current} joint state, calling it during a high-velocity rally
sometimes lands on a different branch than the windup call landed on.
To make consecutive IK calls deterministic we snapshot Robot~1's live
joint positions and velocities, reset the joints to the ready pose for
the IK call, then restore them. No \texttt{stepSimulation} runs in the
intervening window, so the physics state is unchanged.

\subsection{Swing scheduling}
\label{sub:sched}

Given the predicted intercept time $t_{a}=t_{B}+t_{s}$ from the moment
the plan is built, the swing fires
$T_{\mathrm{sw}}/2$ seconds early so that the paddle reaches the strike
point at the midpoint of the windup--follow sweep:
\begin{equation}
  n_{\mathrm{fire}}=\max\!\Bigl(0,\;
   \mathrm{round}\!\bigl((t_{a}-T_{\mathrm{sw}}/2)/h\bigr)\Bigr),
  \qquad h=1/240\,\mathrm{s},\;T_{\mathrm{sw}}=0.22\,\mathrm{s}.
\end{equation}
For $n<n_{\mathrm{fire}}$ the controller commands $\jq_{w}$
(IK of $\mathbf{p}_{w}$); for $n\ge n_{\mathrm{fire}}$ it commands
$\jq_{f}$ (IK of $\mathbf{p}_{f}$). The simulator's position controller
with cap $\dot{\theta}_{\max}=12$\,rad/s then drives the joints between
these targets; with $T_{\mathrm{sw}}=0.22$\,s this comfortably resolves
the largest joint excursion in the plan (about $0.9$\,rad).

\subsection{Net-clearance guard}

The bounce-side test of the original draft was too strict: short returns
landing just past the net (e.g.\ predicted $b_{B}^{x}=+0.03$\,m) would
be rejected even though they are physically returnable. We replaced the
test with a direct \emph{net-clearance check}: if the ball is currently
on the opponent's side, compute the time $t_{N}=-b_{0x}/v_{0x}$ at which
it crosses $x=0$ and the corresponding height
\[
  z_{N}^{\mathrm{ball}}=b_{0z}+v_{0z}\,t_{N}-\tfrac{1}{2}g\,t_{N}^{2}.
\]
Reject the plan only if $z_{N}^{\mathrm{ball}}<0.94$\,m (net top plus a
small ball-radius margin). This single change brought Player~1's plan
acceptance rate from $\sim 4\%$ to a useful fraction without producing
any spurious swings (a ball that fails the clearance test will simply
hit the net regardless of what the arm does).

\subsection{Hyperparameter tuning by exhaustive search}
\label{sub:tune}

The controller exposes eight geometric and timing constants ---
$d_{\mathrm{strike}}$, $\phi$, $d_{W}$, $d_{F}$, $d_{L}$,
$T_{\mathrm{sw}}$, $z^{\min}$ and $e_{B}$ --- whose physically optimal
values depend on quantities (joint controller dynamics, paddle
restitution, ball damping) that are not directly measurable from the
simulator. We sweep all $4\times 3\times 3\times 2\times 3\times 4\times 3\times 2=5184$
combinations on a coarse grid, evaluating each on five seeds in
single-player mode. The sweep is parallelised across 16 worker
processes via \texttt{param\_search.py} and \texttt{param\_search\_runner.py}
and completes in about 16 minutes.

Rather than pick the single best-rate configuration, we choose the
most robust one in the top tier: for each configuration we average its
return rate with that of its grid neighbours (configs differing in
exactly one knob by one step), and select the configuration with the
highest neighbourhood-averaged rate. This guards against picking a
``spike'' that happens to score well on the five seeds but is brittle to
small perturbations.
